\section{Análisis exploratorio de datos}

El catálogo operativo ya está cerrado. La exploración pregunta si la información que la compra necesita está realmente ahí, y qué decisión se sigue de lo que se encuentra. Se usó el archivo oficial y las mismas funciones que usa la aplicación. Cada hallazgo deja una pieza del orden. No es una colección de gráficas sueltas.

\subsection{Cuánta información hay}

La primera pregunta fue cuántos productos se pueden comparar de verdad.

De 13\,093 productos, 7\,568 tienen una categoría de referencia (57{,}80\,\%), esto es, un grupo de productos parecidos contra el cual comparar. 6\,998 tienen al menos cuatro percentiles del núcleo (53{,}45\,\%). Un percentil es el puesto de un nutriente dentro de su categoría, en una escala de 0 a 100. 6\,780 tienen procesamiento calculable (51{,}78\,\%). El universo puntuable junta procesamiento y al menos cuatro de esos percentiles: 5\,864 productos, el 44{,}79\,\% del catálogo. Son productos con los que se puede armar la comparación antes de conocer a la persona. No son productos recomendados. La categoría de referencia no entra en esa regla: es la condición propia de la nutrición. La figura~\ref{fig:disponibilidad} muestra el contraste. Nutrición y procesamiento llegan a cerca de la mitad del catálogo. Etiquetas, alérgenos, trazas y precio cubren mucho menos.

\begin{figure}[ht]
  \centering
  \includegraphics[width=0.92\textwidth]{disponibilidad.png}
  \caption{Productos con dato en cada bloque del catálogo. Fuente: notebook maestro, corte del \corte.}
  \label{fig:disponibilidad}
\end{figure}

De ahí salen tres conjuntos que no se deben mezclar. El catálogo operativo tiene 13\,093 productos. El universo puntuable tiene 5\,864. El orden de una persona es otro conjunto, más chico, y solo aparece cuando ya hay un perfil. Estar en el catálogo no basta para comparar. Estar en el universo puntuable no basta para aparecer en el orden de alguien.

En un corte anterior, del 19 de septiembre de 2026, otra regla más estricta sobre 16\,851 productos dejaba 5\,856 (34{,}8\,\%) en una versión temprana de ese universo. Describe otro archivo. No se usa en las decisiones de este documento.

\subsection{Faltantes}

La cobertura no es pareja. La pregunta siguiente es si esos huecos se pueden reconstruir.

Los huecos de la nutrición, de la categoría y de variables cercanas se concentran en registros que tampoco traen ingredientes, o que Open Food Facts ya marca como incompletos. No aparecen como un valor que se pueda copiar de productos parecidos. Falta información porque el registro está poco documentado, no porque el nutriente valga cero.

Por eso un faltante se guarda como nulo y se marca, para saber que el dato no vino. No se imputan nutrientes, alérgenos, etiquetas, dieta, grupo NOVA ni precio. Imputar sería rellenar el hueco con un estimado. En el orden, una parte sin dato no entra al promedio como cero. Si lo disponible pesa menos de la mitad de lo que la persona pidió, el producto sale de la comparación. Inventar el dato cambiaría el sentido del resultado y, en alergias, el riesgo.

\subsection{Nutrición}

Entre los productos que sí traen nutrientes, no todos sirven para ordenar. Azúcar, sal y grasa saturada tienen una dirección clara dentro de esta comparación: menos es mejor. Fibra y proteína, al revés: más es mejor. Esa dirección solo ordena esta parte del cálculo. No dice que el alimento convenga a cualquier persona.

Energía, grasa total y carbohidratos también tienen percentil y se muestran en la ficha. No hay una meta única que diga si conviene que suban o que bajen para toda la población, así que no se suman. La figura~\ref{fig:nutrientes} muestra la distribución de los cinco que sí entran. La señal nutricional resume esos cinco, no toda la etiqueta.

\begin{figure}[ht]
  \centering
  \includegraphics[width=\textwidth]{nutrientes_d1.png}
  \caption{Distribución de los cinco nutrientes que entran a la dimensión nutricional, en escala de percentil. Fuente: notebook maestro, corte del \corte.}
  \label{fig:nutrientes}
\end{figure}

\subsection{Categorías}

Un nutriente con dirección clara todavía necesita un grupo. Comparar un yogur con todo el catálogo mezclaría alimentos que no compiten en la misma decisión.

Open Food Facts ordena las categorías de lo general a lo específico, por ejemplo de lácteos a un tipo concreto de yogur. Se recorre esa lista desde lo más específico y se para en la primera etiqueta que aparece en al menos 30 productos. En este corte hay 201 categorías de referencia, y ninguna de esas etiquetas está por debajo de 30. Aun así, el grupo con el que se compara un producto puede ser menor, porque no todos los que mencionan la etiqueta la reciben como categoría elegida. La mediana de ese grupo es 27 productos. Hay 5\,525 productos sin categoría de referencia. La figura~\ref{fig:categorias} separa cuántas veces aparece la etiqueta y cuántos productos quedan en el grupo.

\begin{figure}[ht]
  \centering
  \includegraphics[width=0.92\textwidth]{categorias.png}
  \caption{Tamaño del grupo con el que se compara cada producto. Fuente: notebook maestro, corte del \corte.}
  \label{fig:categorias}
\end{figure}

Si ninguna categoría alcanza el mínimo, la nutrición queda sin dato para ese producto. No se compara contra todo el catálogo. El percentil se calcula dentro de la categoría elegida. Sin categoría, la nutrición no aporta al orden.

\subsection{Procesamiento}

La nutrición no agota la comparación. La segunda pieza que la persona puede privilegiar es el grado de procesamiento. NOVA clasifica los alimentos en cuatro grupos, del menos procesado al ultraprocesado \citep{monteiro2018nova}. Aquí esa clasificación ordena una escala. No afirma que un alimento sea adecuado para cualquier persona.

La pregunta era si el grupo, por sí solo, distingue productos en este catálogo. De los que sí tienen grupo, la mayoría está en el 4: 4\,709, frente a 804 en el grupo 1, 189 en el grupo 2 y 1\,078 en el grupo 3. Hay 6\,313 productos sin grupo. Con tanta concentración en un solo grupo, NOVA solo no separa la compra. El número de aditivos, los ingredientes añadidos que el registro enumera, se mueve en el mismo sentido que el grupo y cubre un poco más de productos. La figura~\ref{fig:nova} muestra esa concentración.

\begin{figure}[ht]
  \centering
  \includegraphics[width=0.78\textwidth]{nova.png}
  \caption{Productos por grupo NOVA. Fuente: notebook maestro, corte del \corte.}
  \label{fig:nova}
\end{figure}

El grupo fija un tramo de la escala. Los aditivos mueven al producto dentro de ese tramo, para no dejar juntos a miles de productos del grupo 4. Si no hay grupo, esta parte queda nula: no se inventa un NOVA. Los aditivos no forman una cuarta señal. Solo ajustan el tramo que el grupo ya fijó.

\subsection{Preferencias}

Nutrición y procesamiento se pueden calcular antes de conocer a la persona. Las etiquetas no. Solo tienen sentido frente a lo que esa persona dice que valora.

Solo 3\,776 productos tienen etiquetas. En el perfil de ejemplo, que valora orgánico y sin gluten, 9\,317 productos no tienen con qué comparar y 3\,068 tienen etiquetas, pero ninguna de las dos. Solo 708 presentan al menos una. Si el producto no tiene etiquetas, o si la persona no declaró ninguna, esta parte queda nula. Vale cero solo cuando hay etiquetas y ninguna coincide. No se rellena la preferencia. El cálculo espera al perfil.

\subsection{Precio y calidad}

El precio y la calidad del registro se ven con facilidad, y resulta tentador usarlos para desempatar. Los datos no sostienen ese uso.

El precio de mercado existe en 259 productos (figura~\ref{fig:precios}). Es poco frente al catálogo, así que un orden por precio dejaría fuera a casi todos. La calidad de información, en la figura~\ref{fig:calidad}, describe si el registro está documentado. No mide si el producto encaja con un perfil. Las dos se muestran. No entran al puntaje ni a los filtros del orden. Un precio ausente no se sustituye, dentro del catálogo, por una estimación de mercado.

\begin{figure}[ht]
  \centering
  \includegraphics[width=0.72\textwidth]{precios.png}
  \caption{Precio observado y precio ausente. Fuente: notebook maestro, corte del \corte.}
  \label{fig:precios}
\end{figure}

\begin{figure}[ht]
  \centering
  \includegraphics[width=0.72\textwidth]{calidad.png}
  \caption{Calidad de información del registro. Fuente: notebook maestro, corte del \corte.}
  \label{fig:calidad}
\end{figure}

Cada hallazgo dejó una decisión. La cobertura desigual separa catálogo, universo puntuable y orden de un perfil. Los faltantes se conservan como nulos. La nutrición usa cinco nutrientes y se calcula dentro de la categoría. El procesamiento toma el grupo NOVA y lo ajusta con aditivos, sin inventar el grupo cuando falta. Las etiquetas esperan al perfil. El precio y la calidad se muestran y no ordenan. La preparación de los datos tiene que respetar esas decisiones, no volver a discutirlas.
